% Section 13: open items and next steps
\section{What is open, and what comes next}\label{sec:open}

Sources: PHASE3 \S10, PACKAGE\_DESIGN \S\S6, 8, \path{NEXT_BUILD.md},
\path{EXTERNAL_REVIEW_PLAN.md}, \path{MSIGN_REVIEW_LOG.md}.

\subsection{Open theory}

\paragraph{Conjectures.}
\begin{itemize}[nosep]
\item \textbf{Conjecture P3.5(b3)}: validity of the CLR interval with simulated
candidates (the package default). A reduction to the cited full-candidate CLR
theorem is proved; recovery of rarely optimal vertices, the downward bias of $k_p$
over a subset, growing $|\Lambda|$, and uniform validity are open. Proving it
would make the default a theorem.
\item \textbf{Conjecture N2.2$'$}: the sharp N2 set when size is informative.
\item \textbf{Conjecture C1} (R1): $x^*$ lies in the composition hull ($2\times2$).
\item \textbf{R14}: single-index contextual effects have dimension $C-1$.
\item M1/M2 in regime R3 (high-dimensional Gaussian approximation).
\end{itemize}

\paragraph{Open items from the review rounds.}
\begin{itemize}[nosep]
\item \textbf{R15}: integer unit tables beyond the $2\times2$ margins-only case.
\item \textbf{R16}: sharpness and inference for M-SIGN under district rungs. Until
then grouped M-SIGN is experimental and its intervals have no coverage theorem.
\item Uniform (in the design sequence) limit theory for the plug-in endpoint
beyond P3.4$'$'s localisation condition.
\item Boundary cases where Slater fails: rate and limit of the endpoint estimator;
uniform validity near point identification for CLR (Stoye-type).
\item The rate of the hull error, and inference for the population-hull endpoint.
\item Informative finite-sample inference with many small groups (R2): Theorem P3.8
is valid but uninformative.
\item Calibration of the pooling and M1 tests for a growing number of restrictions
(the wild bootstrap is simulation-supported).
\item Inference for the superpopulation target $B^*$; dependence within districts
beyond the district effect; joint relaxation for population turnover (A1);
coverage of the plug-in under misspecification.
\item The characterisation of strict joint-vs-collapsed gains, and of strict M1
gains, for $R\ge3$.
\end{itemize}

\paragraph{Other proof-register items.} R4 (refutability for $R\ge3$), R5 (general
spanning trees), R6 (strict joint-column gains), R7 (global N2 with estimated
$\tau$), R8 (B1(M), $\kappa$, CAR-$\delta$), R9 (M-SHARP consistency), R11 (DiD).

\subsection{Unreviewed material}

PHASE1b, PHASE2d (review batch~C), PACKAGE\_DESIGN beyond the results used here
and NEW\_AVENUES (batch~D), and a dedicated code review of the M-SIGN and WS5
package changes (batch~E). PHASE2c and PHASE5a were reviewed in Batch~B (two
rounds). The original Phases~1--3 have had two rounds from their referee.

\subsection{Build order (\texttt{NEXT\_BUILD.md})}

\begin{enumerate}[nosep]
\item Release hygiene: a pull request to \texttt{master} and a version tag.
\item M-SHARP, the series-based sharper bound (no new assumption; R9).
\item M-LOC in $2\times2$ and $2\times C$ (needs C1), $2\times C$ joint column
signs, $R\times2$ pairwise ordering on spanning trees.
\item DiD wrappers (gated on review batch~C), L1$'$ quadratic contextual effects,
the non-negative-correlation option for the unit-level term.
\item Later: R16, R15, per-group W2, the drop-one display and the L5/L6 sign check,
the N2 global LP, data-dependent modules (E4/E5, B1(M), $\kappa$, contests,
CAR-$\delta$, dispersion), covariate adjustment, informative R2 inference.
\end{enumerate}
